\section{Conclusion}

Biomedical knowledge graphs and language models are complementary, yet the graph is neither an automatic truth source nor an explanation certificate. The recent literature shows real advances in graph retrieval, community and path reasoning, patient-context prediction, graph-aware representations, model-assisted curation, and verification-first architectures. It also shows that the gains are conditional. Missing or stale knowledge, entity-linking error, poor subgraph selection, lossy rendering, weak evidence use, contradiction, and ungoverned updates can erase the benefit or leave a system less reliable than a simpler alternative.

Progress depends on changing the unit of evaluation from the generated answer to the complete, versioned lifecycle. Studies should establish the graph's causal contribution with matched baselines and ablations, evaluate quality, retrieval, reasoning, generation, operations, and human use separately, preserve provenance and time, and report evidence maturity without equating retrospective benchmarks with clinical validation. Closed-loop systems need quarantined updates, independent verification, reversible versions, and explicit responsibility.

The field should now focus on whether the resulting systems can be inspected, challenged, updated, reproduced, and safely governed. Meeting those criteria would move graph-augmented systems beyond demonstrations and toward trustworthy biomedical AI engineering.
